% Figures: original vector PDFs from the source manuscript, used unchanged.
\begin{figure}[!htbp]
\centering
\includegraphics[width=\linewidth]{figures/figure_1.pdf}
\caption{Study area. (a) Number of tornado and severe-storm major disaster declarations, 1999--2021, designating each county. (b) Small counties, defined as the bottom third of U.S. counties by baseline population. Alaska, Hawaii, and territories are not shown.}
\label{fig:map}
\end{figure}

\begin{figure}[!htbp]
\centering
\includegraphics[width=\linewidth]{figures/figure_2.pdf}
\caption{Participation in HMGP by county population decile. (a) Raw share of county--declarations with any HMGP subapplication and with any funded HMGP project under the declaration. (b) Differences from the largest decile, estimated with declaration fixed effects and hazard and demographic controls (Table~\ref{tab:baseline}, column 4); vertical bars are 95 percent confidence intervals with county-clustered standard errors.}
\label{fig:gradient}
\end{figure}

\clearpage
\begin{figure}[!htbp]
\centering
\includegraphics[width=0.9\linewidth]{figures/figure_3.pdf}
\caption{Where the small-county gap arises. Estimated differences between small and larger counties designated in the same declaration, with full controls including local-government counts (Table~\ref{tab:pipeline}). Circles: application and award outcomes for all county--declarations. Diamonds: shares of decided subapplications approved, withdrawn, or denied, among county--declarations with at least one decided subapplication. Horizontal bars are 95 percent confidence intervals.}
\label{fig:pipeline}
\end{figure}

\begin{figure}[!htbp]
\centering
\includegraphics[width=0.9\linewidth]{figures/figure_4.pdf}
\caption{Within-county response to a tornado hit. Effect of a declaration in which a tornado crossed a populated place in the county, relative to other declarations designating the same county, estimated with county and declaration fixed effects (Equation~\ref{eq:shock}). The small-county response is the sum of the main effect and the interaction, with its standard error computed from the full covariance matrix. Coefficients multiplied by 100; bars are 95 percent confidence intervals.}
\label{fig:shock}
\end{figure}

\clearpage
\begin{figure}[!htbp]
\centering
\includegraphics[width=0.9\linewidth]{figures/figure_5.pdf}
\caption{Small-county gap by year of declaration around the DRRA management-cost provision. Interactions of the small-county indicator with declaration-year indicators (2016 omitted), with county and declaration fixed effects and hazard controls, declarations of 2010--2021. The dashed line marks the date from which DRRA \S1215 reimbursement of subrecipient management costs applies. Bars are 95 percent confidence intervals.}
\label{fig:drra}
\end{figure}
